\section{Data}

\subsection{Target Definition}

v5.2의 corpus universe는 `99` language-script이다. 이 중 `92`개는 XLM-R training language로
표시된 seen/head language이고, `7`개는 XLM-R-unseen target tail language이다. 초기 feedback은
``low-resource language를 task별 3개씩 포함하도록 고른 이유''를 더 명확히 요구했다. 이에
따라 target은 단순히 가장 작은 corpus가 아니라 다음 조건으로 정의한다.

\begin{quote}
XLM-R-unseen tail languages in the smallest Glot500 corpus band that still supports downstream evaluation.
\end{quote}

strict `new\_length <= 10k` non-XLM-R tail language는 local task-list membership이 없어
Tatoeba, Bible, NER, POS 등을 3개씩 채울 수 없다. 따라서 downstream evidence를 확보하기
위해 약 30K--50K Glot500 sentence' band에서 task overlap이 가장 큰 7개 언어를 선택했다.
이 결정은 \codepath{docs/exp/v5.2/1_data_scope/low_resource_task_fill_candidates_ko.md}와
\codepath{preprocessing/prepare_v5_glot50010_merge_inputs.py}의
\codepath{V52_OVERLAP_TASKFILL_GLOT5007}에 기록되어 있다.

\begin{table}[t]
\centering
\caption{Target7 selection. All target languages are XLM-R-unseen Glot500 raw languages with at least 30,000 sentence' units and useful downstream overlap.}
\label{tab:target7}
\resizebox{\textwidth}{!}{%
\begin{tabular}{llllrl}
\toprule
language-script & full name & region & script & new\_length & covered tasks \\
\midrule
dtp\_Latn & Kadazan Dusun & Southeast Asia & Latn & 48,468 & Tatoeba, Bible, Roundtrip, Taxi1500, Embedding similarity \\
xav\_Latn & Xavante & South America & Latn & 31,765 & Bible, Roundtrip, POS, Taxi1500 \\
bam\_Latn & Bambara & West Africa & Latn & 32,150 & Bible, Roundtrip, POS, Taxi1500 \\
csb\_Latn & Kashubian & Europe & Latn & 33,743 & Tatoeba, NER, Embedding similarity \\
ile\_Latn & Interlingue & Constructed/Europe & Latn & 40,984 & Tatoeba, Embedding similarity \\
lij\_Latn & Ligurian & Europe & Latn & 42,447 & NER, POS \\
fur\_Latn & Friulian & Europe & Latn & 30,052 & NER \\
\bottomrule
\end{tabular}}
\end{table}


\subsection{Task Coverage}

Target7은 총 21개 task slot을 7개 언어로 압축한다. 이 표는 왜 downstream task를 사용할 수
있는지, 그리고 어떤 결과를 target evidence로 말할 수 있는지의 근거이다.

\begin{table}[t]
\centering
\caption{Downstream coverage rationale for Target7. This table is the defense artifact for why the target set is not the strict smallest corpus set.}
\label{tab:task-coverage}
\begin{tabularx}{\textwidth}{lXl}
\toprule
Task & target languages & current status \\
\midrule
Tatoeba retrieval & csb\_Latn, dtp\_Latn, ile\_Latn & ready local pairs \\
Bible retrieval & dtp\_Latn, xav\_Latn, bam\_Latn & materialized parallel verses \\
Roundtrip alignment & dtp\_Latn, xav\_Latn, bam\_Latn & Bible-derived JSONL \\
NER & csb\_Latn, lij\_Latn, fur\_Latn & PAN-X train/dev/test ready \\
POS & xav\_Latn, bam\_Latn, lij\_Latn & rebuilt/materialized split \\
Taxi1500 & dtp\_Latn, xav\_Latn, bam\_Latn & not target7 evidence in current draft \\
Embedding similarity & csb\_Latn, dtp\_Latn, ile\_Latn & diagnostic pairs ready \\
\bottomrule
\end{tabularx}
\end{table}


\subsection{Training Corpus}

MLM training corpus는 \codepath{preprocessing/run_v52_glot5007_merge.sh}가
\codepath{preprocessing/merge_files.py}를 호출하여 만든다. Sampling은 Glot500-style language
sampling factor $\alpha=0.3$을 사용한다.

\[
P_i \propto n_i^\alpha
\]

여기서 $n_i$는 language-script $i$의 corpus size이다. $\alpha=1$이면 head language가 corpus
크기만큼 압도적으로 자주 뽑히지만, $\alpha=0.3$은 head/tail ratio를 완화한다. 예를 들어
1,000,000 sentence head와 1,000 sentence tail은 $\alpha=1$에서 1000:1이지만, $\alpha=0.3$에서는
약 8:1 수준으로 줄어든다. v5.2 merge report의 핵심 수치는 다음과 같다.

\begin{center}
\begin{tabular}{lr}
\toprule
Item & Samples \\
\midrule
seen/head training samples & 4,147,176 \\
target7 training samples & 335,083 \\
total MLM training samples & 4,482,259 \\
\bottomrule
\end{tabular}
\end{center}

이 값은 \codepath{docs/exp/v5.2/LIVE_STATUS.md} 및
\codepath{docs/exp/v5.2/0_tokenizer/merge/Glot500_v52_glot5007_xlmr100.report.json}의
ready/PASS 상태와 일치한다.
